% sections/06_overfitting_control.tex  –  Overfitting Control
% \input{} into main.tex; do NOT wrap in \documentclass or \begin{document}

\section{Overfitting Control and Evaluation Integrity}
\label{sec:overfit}

With only \num{1603} labelled traces, this project operates in a regime
where overfitting and evaluation bias are the \emph{primary} threats to
credible conclusions.  A model that memorises the training set —
or a cross-validation procedure that leaks information across folds —
can produce accuracy figures that bear no relationship to real-world
performance on the edge sensor.  This section catalogues the safeguards
that the experimental pipeline must enforce.

%% ------------------------------------------------------------------
\subsection{Context: Small-Sample Deep Learning}
\label{sec:overfit:context}

Modern deep architectures routinely contain millions of trainable
parameters.  A ResNet-1D or 1D~CNN with even a modest channel depth can
easily exceed $10^{6}$ weights, yielding a parameter-to-sample ratio far
above unity.  Classical statistical learning theory warns that
generalisation degrades when model capacity greatly exceeds the
information content of the training set~\cite{vapnik1995svm}.  While
modern regularisation techniques mitigate this to a degree, they do not
eliminate the risk.

The practical consequence is twofold:
\begin{enumerate}
  \item \textbf{Architecture selection} should start from the smallest
        viable model and scale up only if validation performance
        saturates (\Cref{sec:overfit:capacity}).
  \item \textbf{Evaluation protocols} must be unusually rigorous; a
        single random train/test split is insufficient for stable
        performance estimates when $N = 1603$
        (\Cref{sec:overfit:cv}).
\end{enumerate}

%% ------------------------------------------------------------------
\subsection{Stratified Train / Validation / Test Split}
\label{sec:overfit:split}

The simplest holdout protocol divides the dataset into three disjoint
partitions while preserving the class proportions
($601:501:501 \approx 1.20{:}1{:}1$):

\begin{center}
\begin{tabular}{lccc}
\toprule
\textbf{Partition} & \textbf{Share} & \textbf{Purpose} \\
\midrule
Training      & $60$--$70\,\%$ & Fit model parameters \\
Validation    & $15$--$20\,\%$ & Tune hyper-parameters, early stopping \\
Test          & $15$--$20\,\%$ & Final, single-use performance estimate \\
\bottomrule
\end{tabular}
\end{center}

Stratification ensures that every partition contains approximately the
same proportion of \texttt{dry}, \texttt{wet}, and \texttt{ice} samples.
For $N=1603$ and a $60/20/20$ split, the test set contains roughly
$320$ samples — enough for per-class precision and recall estimates but
small enough that a single unlucky split can shift accuracy by several
percentage points.

%% ------------------------------------------------------------------
\subsection{Repeated Stratified $k$-Fold Cross-Validation}
\label{sec:overfit:cv}

To obtain more stable performance estimates, we employ \emph{repeated
stratified $k$-fold cross-validation}~\cite{kohavi1995cv}.  Two
configurations are considered:

\begin{itemize}
  \item \textbf{$5 \times 5$-fold:} five repetitions of $5$-fold CV,
        yielding $25$ performance measurements.  Each fold contains
        ${\sim}320$ test samples.
  \item \textbf{$3 \times 10$-fold:} three repetitions of $10$-fold CV,
        yielding $30$ measurements with ${\sim}160$ test samples per
        fold.
\end{itemize}

The mean and standard deviation (or $95\,\%$ confidence interval) across
all folds are reported.  A high standard deviation signals that the
model's performance is sensitive to the particular train/test partition,
which is itself a diagnostic of overfitting or insufficient data.

\paragraph{Critical implementation detail.}
All preprocessing that involves fitting to data — normalization
statistics (\Cref{sec:norm}), feature-extraction parameters, and
dimensionality-reduction transforms — must be fit \emph{exclusively} on
the training fold of each CV iteration and re-fit from scratch when the
fold changes.  Fitting once on the full dataset and then splitting
constitutes data leakage and invalidates the CV estimate.

%% ------------------------------------------------------------------
\subsection{Nested Cross-Validation for Hyper-Parameter Tuning}
\label{sec:overfit:nested}

When hyper-parameters (e.g., regularisation strength, learning rate,
tree depth) are tuned via grid or random search, an additional inner CV
loop is required to prevent the test-fold performance from being
confounded with the tuning process.  This is \emph{nested
cross-validation}~\cite{varma2006nested_cv,cawley2010overfitting}:

\begin{enumerate}
  \item \textbf{Outer loop} ($k_{\text{out}}$ folds): reserves one fold
        as the test set.
  \item \textbf{Inner loop} ($k_{\text{in}}$ folds on the remaining
        data): selects the best hyper-parameter configuration by
        validation performance.
  \item The best configuration from the inner loop is retrained on the
        full outer-training set and evaluated on the outer test fold.
\end{enumerate}

A typical configuration is $k_{\text{out}} = 5$,
$k_{\text{in}} = 3$, yielding $15$ model fits per hyper-parameter
candidate.  The computational cost is substantial but necessary:
Cawley and Talbot~\cite{cawley2010overfitting} demonstrated that
non-nested CV can overfit the \emph{evaluation metric itself},
producing optimistic estimates that do not transfer to unseen data.

%% ------------------------------------------------------------------
\subsection{Near-Duplicate and Temporal-Correlation Awareness}
\label{sec:overfit:duplicates}

The trace file names (e.g., \texttt{Trc100.txt}, \texttt{Trc101.txt})
suggest sequential acquisition.  If consecutive traces were recorded
within seconds of each other under identical sensor conditions, they may
be near-duplicates: statistically indistinguishable samples that, if
split across training and test partitions, artificially inflate
performance.

\paragraph{Detection protocol.}
\begin{enumerate}
  \item Compute the pairwise Pearson correlation (or cosine similarity)
        of the magnitude vectors within each class.
  \item Identify clusters of traces with inter-trace correlation
        $r > 0.999$ (or another empirically determined threshold).
  \item Assign all members of a high-correlation cluster to the
        \emph{same} CV fold using a \emph{group-aware} splitter (e.g.,
        \texttt{GroupKFold} in scikit-learn~\cite{sklearn_docs}).
\end{enumerate}

Group-aware splitting ensures that the model cannot trivially recognise
a test trace by having seen its near-twin during training.  The trace-
numbering gaps noted in \Cref{sec:preproc} (\texttt{Trc601}--\texttt{Trc699},
\texttt{Trc1201}--\texttt{Trc1449}) may further indicate batch
boundaries that should be respected during splitting.

%% ------------------------------------------------------------------
\subsection{Capacity--Dataset Size Mismatch}
\label{sec:overfit:capacity}

\Cref{tab:model_capacity} provides rough parameter counts for the
candidate architectures.  The table is indicative; exact counts depend
on final architecture choices.

\begin{table}[ht]
\centering
\caption{Indicative parameter counts and overfitting risk relative to
         $N=1603$ samples.}
\label{tab:model_capacity}
\small
\begin{tabular}{lrll}
\toprule
\textbf{Algorithm} &
\textbf{Approx.\ parameters} &
\textbf{Param./sample} &
\textbf{Overfit risk} \\
\midrule
Logistic Regression     & ${\sim}\num{12000}$   & ${\sim}7$     & Low      \\
Linear SVM              & ${\sim}\num{12000}$   & ${\sim}7$     & Low      \\
RBF SVM                 & support vectors       & —             & Moderate \\
Random Forest           & ensemble, depth-bound  & —             & Low--Mod \\
Gradient Boosted Trees  & ensemble, depth-bound  & —             & Moderate \\
Multilayer Perceptron   & $\num{50000}$--$\num{500000}$  & $30$--$300$  & High \\
1D CNN                  & $\num{100000}$--$\num{2000000}$ & $60$--$1200$ & High \\
ResNet-1D               & $\num{500000}$--$\num{5000000}$ & $300$--$3000$ & Very high \\
KAN                     & architecture-dependent & variable      & High \\
\bottomrule
\end{tabular}
\end{table}

\paragraph{Recommendation.}
Begin every deep-learning experiment with the \emph{smallest} viable
architecture (e.g., two convolutional layers with $16$--$32$ filters for
the 1D~CNN).  Scale up only if validation loss plateaus without
overfitting.  Tree-based ensembles and linear models serve as strong
baselines precisely because their effective capacity is easier to
control.

%% ------------------------------------------------------------------
\subsection{Regularization Toolkit}
\label{sec:overfit:reg}

The following regularisation techniques are available and should be
applied in combination as needed.

\subsubsection{Early Stopping}
\label{sec:overfit:reg:es}

Monitor the validation loss at the end of each training epoch.  If the
loss does not improve for a predefined number of consecutive epochs
(the \emph{patience}), training is halted and the model weights from
the best epoch are restored.  A patience of $10$--$20$ epochs is a
reasonable starting point for this dataset size; excessively large
patience allows the model to overfit before the stopping criterion
triggers.

\subsubsection{Weight Decay / L2 Regularization}
\label{sec:overfit:reg:wd}

Weight decay adds a penalty proportional to the squared $\ell_2$-norm
of the model weights to the loss function:
\begin{equation}
  \mathcal{L}_{\text{reg}} = \mathcal{L}_{\text{task}}
      + \frac{\lambda}{2}\|\mathbf{w}\|_2^2.
  \label{eq:weight_decay}
\end{equation}
For adaptive optimisers (Adam, AdamW), decoupled weight
decay~\cite{loshchilov2017decoupled} is preferred because it disentangles
the regularisation effect from the adaptive learning rate.

\subsubsection{Dropout}
\label{sec:overfit:reg:dropout}

Dropout~\cite{srivastava2014dropout} randomly zeroes a fraction $p$ of
activations during each forward pass, forcing the network to learn
redundant representations.  Typical values are $p \in [0.2,\,0.5]$.
For very small datasets, moderate dropout ($p \approx 0.3$) combined
with early stopping often outperforms aggressive dropout, which can
under-fit.

\subsubsection{Label Smoothing}
\label{sec:overfit:reg:ls}

Standard one-hot targets encourage the model to produce overconfident
predictions (logit magnitudes $\to \infty$).  Label smoothing replaces
the hard target $y_c \in \{0,1\}$ with a soft target:
\begin{equation}
  y_c^{\text{smooth}} = (1 - \alpha)\,y_c + \frac{\alpha}{C},
  \label{eq:label_smooth}
\end{equation}
where $C=3$ is the number of classes and $\alpha \in [0.05,\,0.2]$ is
the smoothing factor.  This reduces the gradient magnitude for
already-correct predictions and acts as an implicit regulariser.

\subsubsection{Pruning}
\label{sec:overfit:reg:pruning}

Post-training weight pruning removes parameters that contribute little
to the output, simultaneously reducing model size and (potentially)
improving generalisation by eliminating memorised noise.  Pruning is
especially relevant for edge deployment (\Cref{sec:nextsteps}), where
model size and inference latency are constrained.

%% ------------------------------------------------------------------
\subsection{Confidence Calibration as a Diagnostic}
\label{sec:overfit:calibration}

A well-calibrated model produces predicted probabilities that match
empirical class frequencies: if the model assigns $p=0.8$ to a class
across many test samples, roughly $80\,\%$ of those predictions should
be correct.  Guo~et~al.~\cite{guo2017calibration} showed that modern
deep networks are frequently \emph{miscalibrated} — typically
overconfident — even when their accuracy is high.

Calibration diagnostics (reliability diagrams, Expected Calibration
Error) are \textbf{useful post-training assessments} of prediction
quality but are \textbf{not substitutes} for proper generalisation
control.  A model can be well-calibrated yet still overfit to a leaky
evaluation protocol.  Calibration should therefore be measured
\emph{after} all the safeguards in this section are in place, using the
held-out test set or the outer fold of nested CV.

Post-hoc calibration methods (temperature scaling, Platt scaling) can
correct miscalibration at minimal computational cost and are recommended
before edge deployment, where the classifier's confidence score may
influence downstream decision logic (e.g., ``flag for human review if
$\max p_c < 0.7$'').

%% ------------------------------------------------------------------
\subsection{Summary of Safeguards}
\label{sec:overfit:summary}

\begin{tcolorbox}[
  colback=yellow!5!white,
  colframe=yellow!70!black,
  title={\textbf{Key Principle}},
  fonttitle=\bfseries
]
Small-sample edge-sensing work is \emph{especially} vulnerable to
optimistic performance estimates produced by careless data splits,
leaked preprocessing, or unchecked model capacity.  Every reported
metric in this project must survive the following audit:
\begin{enumerate}[nosep]
  \item Was normalization fit only on the training fold?
  \item Were near-duplicate traces kept in the same fold?
  \item Was hyper-parameter tuning isolated in an inner CV loop?
  \item Is the reported metric averaged over multiple folds/repeats?
  \item Is the model's parameter count justified relative to $N=1603$?
\end{enumerate}
A negative answer to any of these questions invalidates the associated
result.
\end{tcolorbox}
